%!TEX program = uplatex
\RequirePackage{plautopatch}
\documentclass[uplatex,dvipdfmx]{jlreq}
\usepackage{amsmath,amssymb}

\pagestyle{empty}

\begin{document}

%======================================================================
% 表紙
%======================================================================
\begin{center}
  \vspace*{20pt}
  {\Huge\bfseries 不変式論のルネサンス}

  \vspace{1pc}
  {\Large 概要}
\end{center}

\vfill

\newpage

%======================================================================
% 講演１：梅田亨「不変式論の概観 --- 歴史と枠組み」
%======================================================================
\begin{center}
  {\Large\bfseries 不変式論の概観 --- 歴史と枠組み}
\end{center}

\bigskip
\hspace*{\fill}{\large\bfseries 梅田　亨（京大・理）}

\bigskip

（代数的）不変式論は今から丁度160年前、G.~Boole が
$m$ 元 $n$ 次形式の判別式を定義し、その（相対）不変性を示した1841年に始まる。
A.~Cayley はその論文の重要性を認め、直ちに不変式の組織的研究を始めた。
同時にそれは時代の流れに乗って多くの数学者の関心を惹き、非常に多くの研究がなされた。
初期の結果で重要なのは、いわゆる記号的方法で、
それによって原理的には全ての不変式を書き下すことが可能となった。
ついで問題は不変式をより効率的に書くこと、特に環論的な意味での生成系を知ること、
に移り、その有限性が焦点となった。
P.~Gordan は不変式の王と言われ、特に二元形式（binary form）の不変式の生成系の
有限性を示した。
それは構成的な方法であったが、技術的な限界にも近づいていた。
そのような中 D.~Hilbert が多項式イデアルに関する基底定理（1888）を以て、
全く一般的に問題を解決してしまった。
これは「神学」と呼ばれたほどの抽象的性格を持ち、非構成的なものである。
この新しい視点は現代数学の幕開けをもたらした。

\medskip

Hilbert はさらに零点定理（Nullstellensatz）や syzygy 定理などを以て、
不変式論を深めた
（例えば零点定理は、非構成的定理から構成的定理への橋渡しである）が、
そのようにして主要な問題が解かれてしまったことで、
不変式論の研究は急速に下火になった。
いわゆる「不変式論の死」である。
しかし、実際は不変式論は死んではおらず、現代に至るまで何度も再生を繰り返している。
今回のテーマ「不変式論のルネサンス」はそのような文脈からつけられた。
そしてその一端の幾つかは今回の講演で実際触れられることになる。

\medskip

ここでは、まず全体の前座として大雑把な概観を見たい。
不変式論については、共通の知識を仮定できるほど標準的なものとなっていない現状を踏まえ、
歴史を若干溯ることと、現代数学の標準的な言葉での枠組みを述べることで、
予備的導入としたい。

\newpage

%======================================================================
% 講演２：向井茂「不変式の名所旧跡」
%======================================================================
\begin{center}
  {\Large\bfseries 不変式の名所旧跡}
\end{center}

\bigskip
\hspace*{\fill}{\large\bfseries 向井　茂}

\bigskip

不変式の世界は広くて整理が大変ですが、群から見ると次の３つに大別されます。
\begin{enumerate}
  \item 対称式や交代式のように有限群で不変なもの。
  \item 行列式や判別式のように連続群（たいていは座標変換）で不変なもの。
  \item $\wp$ 関数やテータのように無限離散群で不変なもの。
\end{enumerate}

1で訪れる人の多いのは鏡映群に関するものですが、Weyl 群が鏡映群なので2と絡み合います。
正多面体群も有名です。これの中心拡大が平面に作用して、重要な特異点を作りだします。
これの発展として3次特殊線型群 $SL(3)$ の有限部分群の不変式も注目を浴びていますが、
ここには位数168の単純群で不変な Klein の4次式
\[
  x^3 y + y^3 z + z^3 x
\]
という別の名所があります。

1、2、3のいずれにも Heisenberg 群という特徴的で重要な群がありますが、
3変数3次式
\[
  x^3 + y^3 + z^3 - 3\lambda xyz, \qquad \lambda : \text{定数}
\]
や4変数4次式
\[
  x^4 + y^4 + z^4 + t^4 + A(x^2 y^2 + z^2 t^2)
  + B(x^2 z^2 + y^2 t^2) + C(x^2 t^2 + y^2 z^2) + Dxyzt
\]
が有限 Heisenberg 群の不変式として特に有名です。

2は他の講演で沢山出てきますので、基本だけを押さえましょう。
少しは、概均質ベクトル空間に触れたいと思います。

3はほとんど保型関数論です。
$\wp$ 関数は平行移動群 $\mathbf{Z}^2$ で不変な有理型関数で、
$\theta(\tau, z)$ は Heisenberg 群（$\mathbf{Z}^2$ の中心拡大）の作用で不変な正則関数です。
また、その零値 $\theta(\tau, 0)$ は $SL(2, \mathbf{Z})$ の適当な部分群に対して
ある不変性をもちます。

この Encounter with Math.~は2が中心になりそうですが、
それ以外の所もぜひ訪れてみて下さい。

\newpage

%======================================================================
% 講演３：寺西鎮男「"The classical groups" を逆に読む」
%======================================================================
\begin{center}
  {\Large\bfseries ``The classical groups'' を逆に読む}
\end{center}

\bigskip
\hspace*{\fill}{\large\bfseries 寺西　鎮男（名古屋大・多元数理）}

\bigskip

Herman Weyl によれば、不変式論は Hilbert による有限性定理以来、何度も死を宣告されながら
幾度も再生してきた。
ある数学的な対象に対して、様々な見方ができることが興味ある数学的対象であるための
一つの条件であるとすれば、不変式論には少なくとも次の4つの側面がある：
\begin{enumerate}
  \item 組み合わせ論的側面
  \item 幾何学的側面
  \item 環論的側面
  \item 群の表現論的側面
\end{enumerate}

初期の段階においては、組み合わせ論的側面が主として研究され
（Cayley、Sylvester による記号的方法、不変式の個数定理等）、
それが群の表現論の発展を促した。
Hilbert は彼の基底定理を用いて当時の懸案の問題であった不変式の有限性定理を一般的に証明し、
更に零点定理等を用いて不変式論の環論的側面や幾何学的側面を開拓した。

Weyl は \textit{The classical groups}~[W] において古典群に対する、
ベクトル不変式の基本定理を Capelli の恒等式を用いて証明し、
それを古典群の表現論に応用した。
例えば、一般線形群の表現と対称群の表現は、Schur の相互律により結びつけられるが、
Schur の相互律は一般線形群のベクトル不変式の第一基本定理に他ならない。

この講演では、Weyl の本を逆に進んで、
群の表現論から古典的な不変式論の基本的な諸結果を見直すことを試みる。
より具体的には表現論の基本的な結果を説明したのち、次の話題にふれる予定である。
\begin{itemize}
  \item binary form の不変式
  \item semi-invariant と既約表現
  \item ベクトル不変式の基本定理と記号的方法
  \item matrix invariant と永田--Higman index
\end{itemize}

\newpage

%======================================================================
% 講演４：梅田亨「不変式と双対性」
%======================================================================
\begin{center}
  {\Large\bfseries 不変式と双対性}
\end{center}

\bigskip
\hspace*{\fill}{\large\bfseries 梅田　亨（京大・理）}

\bigskip

H.~Weyl はその古典的書物
\textit{The Classical Groups --- their invariants and representations}（1939）において
不変式論の再生を試みた。
この書物の元になった Princeton 講義録（1936/37）のタイトル
\textit{Elementary Theory of Invariants} がその意図をより明瞭に表わしているが、
表現論という基本的な枠組みを打ちたてつつ、
不変式論と表現論の相互のつながりを密接な形に描き出している。
不変式論から言えば、特に基本的な対象であるベクトル不変式についての
諸結果（特に第一及び第二基本定理）が、典型群（classical groups）に対して確立されている。

\medskip

これらの結果は、しかしながら、使われた道具が幾分形式的な Capelli Identity というものであった
せいもあって、その後しばらくは蝕にあっていたようである。
それらが更にはっきりした表現論的意味と深い応用を伴って我々の前に登場したのは
R.~Howe による Dual pair 理論（1976）の文脈に於いてである。
それは保型形式論や物理学に現われるさまざまな双対性という現象を統一的に捉える視点を与えるが、
その元型が Weyl の本で確立されている第一基本定理だったのである。

\medskip

本講演では、この不変式論の第一基本定理と、その双対定理としての側面について解説したい。

\newpage

%======================================================================
% 講演５：向井茂「不変式環 --- Hilbert の例と永田の例」
%======================================================================
\begin{center}
  {\Large\bfseries 不変式環 --- Hilbert の例と永田の例}
\end{center}

\bigskip
\hspace*{\fill}{\large\bfseries 向井　茂}

\bigskip

Hilbert が不変式論を研究したのは32歳までの約8年間ですが、その影響は甚大でした。
ここでは論文 [Hi] と彼の提出した数学の23問題の一つ

\medskip
\begin{center}
  「代数群が多項式環に作用するとき、その不変式環は有限生成か？」
\end{center}
\medskip

\noindent
について考えます。

Hilbert の視野には代数関数論の多変数への拡張があり、不変式体は
「数体に対する重要な命題が円分体において最初に発見され、そして、証明されるのと
同じような意味で重要な」多変数関数体の例であると序文で述べています。
代数多様体について多くのことが分かった現在においても味わいある言葉でしょう。
この論文の第4・5章は Mumford によって安定性とその数値的判定法に一般化されています。
第3章は、基本的な不変式環に対して、その Hilbert 級数と重複度
（商多様体の次数といっても同じ）を計算しています。

さて、Hilbert の提出した上の問題は永田 [Na] によって否定的な例が与えられました。
これはいろんな数学と（不思議なことに [Hi] 第3章とも）関連する健康な例のようです。
[Mu] 第2章をさらに簡略化した形で紹介します。

\newpage

%======================================================================
% 参考文献（共通）
%======================================================================
\begin{center}
  {\large\bfseries 不変式論のルネサンス}
\end{center}

\begin{center}
  \textsc{参考文献}
\end{center}

\begin{enumerate}
  \item[{[GY]}] I.~H.~Grace and A.~Young,
        \textit{The Algebra of Invariants},
        AMS/Chelsea, 1903.

  \item[{[Ha]}] T.~Hawkins,
        \textit{Emergence of the Theory of Lie Groups:
        an essay in the history of Mathematics 1869--1926},
        Springer.

  \item[{[Hi]}] D.~Hilbert,
        Über die vollen Invariantensysteme,
        \textit{Math.~Ann.}~\textbf{42} (1893), 313--373.

  \item[{[Ho1]}] R.~Howe,
        Remarks on classical invariant theory,
        \textit{Trans.~Amer.~Math.~Soc.}~\textbf{313} (1989), 539--570.
        Erratum: \textit{ibid.}~\textbf{318} (1990), 823.

  \item[{[Ho2]}] R.~Howe,
        Perspectives on invariant theory: Schur duality,
        multiplicity-free actions and beyond,
        in \textit{The Schur Lectures (1992)},
        Israel Math.~Conf.~Proc., vol.~8,
        Bar-Ilan University, 1995, 1--182.

  \item[{[Mu]}] 向井茂：
        モジュライ理論 1, 2, 岩波書店, 東京, 1998年・2000年.

  \item[{[Na]}] M.~Nagata,
        On the fourteenth problem of Hilbert,
        \textit{Int'l~Cong.~Math.}, Edinburgh, 1958.

  \item[{[R]}] C.~リード：
        ヒルベルト --- 現代数学の巨峰 ---（弥永健一訳）、岩波書店、1972年.

  \item[{[U]}] 梅田亨：
        不変式論・入門・以前 --- 第１基本定理と記号的方法 ---、
        論集　現代の母関数、1991年、pp.~71--188.

  \item[{[W]}] H.~Weyl,
        \textit{The Classical Groups: Their Invariants and Representations},
        Princeton University Press, Princeton, 1946.
\end{enumerate}

\end{document}
